% Target: rmp-ii-peps-marginal
% Formula: marginals rho_kept = Tr_traced |psi><psi|.
% Ink: Kernel plane frames; the doubled sheet is the contract bilayer basis
%   (planes preset) with member-addressed topology.
% Boundary: The state sheet is open; on the doubled sheet the kept sites'
%   ket-bra physical legs stay open and the fourth column closes site by site
%   through east wrap loops.
% Source: paper TN-Review-main.tex line 331; author source ImagesReview Section
%   II/ImagesReview Section II.tex lines 287-304
% Capabilities: kernel, plane-frame
\tenkzkernel \[ \begin{gathered}
  % The state sheet |psi>: every site keeps its transverse physical leg.
  \begin{tenkz}[lattice={4x4}, frame=plane, physical=up, boundary=open]
    \tn{} & \tn{} & \tn{} & \tn{} \\
    \tn{} & \tn{} & \tn{} & \tn{} \\
    \tn{} & \tn{} & \tn{} & \tn{} \\
    \tn{} & \tn{} & \tn{} & \tn{}
  \end{tenkz} \\[1em]
  % The doubled sheet |psi><psi|: the plane frame owns both sheets via the
  % bilayer basis; every bond, stub, leg, and traced pair below is member-addressed.
  \begin{tenkz}[lattice={4x4}, planes, bonds=none]
    % ket-sheet east-west bonds
    \tnwire{(1,1,1)}{(1,2,1)} \tnwire{(1,2,1)}{(1,3,1)} \tnwire{(1,3,1)}{(1,4,1)}
    \tnwire{(2,1,1)}{(2,2,1)} \tnwire{(2,2,1)}{(2,3,1)} \tnwire{(2,3,1)}{(2,4,1)}
    \tnwire{(3,1,1)}{(3,2,1)} \tnwire{(3,2,1)}{(3,3,1)} \tnwire{(3,3,1)}{(3,4,1)}
    \tnwire{(4,1,1)}{(4,2,1)} \tnwire{(4,2,1)}{(4,3,1)} \tnwire{(4,3,1)}{(4,4,1)}
    % ket-sheet north-south bonds
    \tnwire{(1,1,1)}{(2,1,1)} \tnwire{(2,1,1)}{(3,1,1)} \tnwire{(3,1,1)}{(4,1,1)}
    \tnwire{(1,2,1)}{(2,2,1)} \tnwire{(2,2,1)}{(3,2,1)} \tnwire{(3,2,1)}{(4,2,1)}
    \tnwire{(1,3,1)}{(2,3,1)} \tnwire{(2,3,1)}{(3,3,1)} \tnwire{(3,3,1)}{(4,3,1)}
    \tnwire{(1,4,1)}{(2,4,1)} \tnwire{(2,4,1)}{(3,4,1)} \tnwire{(3,4,1)}{(4,4,1)}
    % bra-sheet east-west bonds
    \tnwire{(1,1,2)}{(1,2,2)} \tnwire{(1,2,2)}{(1,3,2)} \tnwire{(1,3,2)}{(1,4,2)}
    \tnwire{(2,1,2)}{(2,2,2)} \tnwire{(2,2,2)}{(2,3,2)} \tnwire{(2,3,2)}{(2,4,2)}
    \tnwire{(3,1,2)}{(3,2,2)} \tnwire{(3,2,2)}{(3,3,2)} \tnwire{(3,3,2)}{(3,4,2)}
    \tnwire{(4,1,2)}{(4,2,2)} \tnwire{(4,2,2)}{(4,3,2)} \tnwire{(4,3,2)}{(4,4,2)}
    % bra-sheet north-south bonds
    \tnwire{(1,1,2)}{(2,1,2)} \tnwire{(2,1,2)}{(3,1,2)} \tnwire{(3,1,2)}{(4,1,2)}
    \tnwire{(1,2,2)}{(2,2,2)} \tnwire{(2,2,2)}{(3,2,2)} \tnwire{(3,2,2)}{(4,2,2)}
    \tnwire{(1,3,2)}{(2,3,2)} \tnwire{(2,3,2)}{(3,3,2)} \tnwire{(3,3,2)}{(4,3,2)}
    \tnwire{(1,4,2)}{(2,4,2)} \tnwire{(2,4,2)}{(3,4,2)} \tnwire{(3,4,2)}{(4,4,2)}
    % ket-sheet perimeter openings
    \tnwire{(1,1,1)}{open n} \tnwire{(1,2,1)}{open n} \tnwire{(1,3,1)}{open n}
    \tnwire{(1,4,1)}{open n} \tnwire{(4,1,1)}{open s} \tnwire{(4,2,1)}{open s}
    \tnwire{(4,3,1)}{open s} \tnwire{(4,4,1)}{open s} \tnwire{(1,1,1)}{open w}
    \tnwire{(2,1,1)}{open w} \tnwire{(3,1,1)}{open w} \tnwire{(4,1,1)}{open w}
    \tnwire{(1,4,1)}{open e} \tnwire{(2,4,1)}{open e} \tnwire{(3,4,1)}{open e}
    \tnwire{(4,4,1)}{open e}
    % bra-sheet perimeter openings
    \tnwire{(1,1,2)}{open n} \tnwire{(1,2,2)}{open n} \tnwire{(1,3,2)}{open n}
    \tnwire{(1,4,2)}{open n} \tnwire{(4,1,2)}{open s} \tnwire{(4,2,2)}{open s}
    \tnwire{(4,3,2)}{open s} \tnwire{(4,4,2)}{open s} \tnwire{(1,1,2)}{open w}
    \tnwire{(2,1,2)}{open w} \tnwire{(3,1,2)}{open w} \tnwire{(4,1,2)}{open w}
    \tnwire{(1,4,2)}{open e} \tnwire{(2,4,2)}{open e} \tnwire{(3,4,2)}{open e}
    \tnwire{(4,4,2)}{open e}
    % the kept sites' 24 open ket-bra physical legs: rho's row indices up
    % from the ket sheet, its column indices down from the bra sheet
    \tnwire{(1,1,1)}{open up} \tnwire{(1,2,1)}{open up} \tnwire{(1,3,1)}{open up}
    \tnwire{(2,1,1)}{open up} \tnwire{(2,2,1)}{open up} \tnwire{(2,3,1)}{open up}
    \tnwire{(3,1,1)}{open up} \tnwire{(3,2,1)}{open up} \tnwire{(3,3,1)}{open up}
    \tnwire{(4,1,1)}{open up} \tnwire{(4,2,1)}{open up} \tnwire{(4,3,1)}{open up}
    \tnwire{(1,1,2)}{open down} \tnwire{(1,2,2)}{open down} \tnwire{(1,3,2)}{open down}
    \tnwire{(2,1,2)}{open down} \tnwire{(2,2,2)}{open down} \tnwire{(2,3,2)}{open down}
    \tnwire{(3,1,2)}{open down} \tnwire{(3,2,2)}{open down} \tnwire{(3,3,2)}{open down}
    \tnwire{(4,1,2)}{open down} \tnwire{(4,2,2)}{open down} \tnwire{(4,3,2)}{open down}
    % The traced fourth column: each ket-bra pair closes through its own
    % east wrap loop (Tr over the site's physical space), restored here as
    % the paper's three-sided wrap: two invisible junction atoms per loop
    % chained by three kind=string segments (tN, tNb, tNc).  Consecutive
    % segments meet at a shared junction, which the crossing police reads
    % as an undeclared crossing, so each adjacent pair carries a nominal
    % cross= declaration; the outer segments (tN, tNc) go under so the gap
    % trim lands on them, not on the short middle jump tNb.
    \tn[at=0.6 e of (1,4,1), skin=none, name=j11]{}
    \tn[at=0.6 e of (1,4,2), skin=none, name=j12]{}
    \tnwire[kind=string, name=t1, cross={under at crossing of t1 and t1b}]{(1,4,1)}{j11}
    \tnwire[kind=string, name=t1b]{j11}{j12}
    \tnwire[kind=string, name=t1c,
            cross={under at crossing of t1c and t1b}]{j12}{(1,4,2)}
    \tn[at=0.6 e of (2,4,1), skin=none, name=j21]{}
    \tn[at=0.6 e of (2,4,2), skin=none, name=j22]{}
    \tnwire[kind=string, name=t2, cross={under at crossing of t2 and t2b}]{(2,4,1)}{j21}
    \tnwire[kind=string, name=t2b]{j21}{j22}
    \tnwire[kind=string, name=t2c,
            cross={under at crossing of t2c and t2b}]{j22}{(2,4,2)}
    \tn[at=0.6 e of (3,4,1), skin=none, name=j31]{}
    \tn[at=0.6 e of (3,4,2), skin=none, name=j32]{}
    \tnwire[kind=string, name=t3, cross={under at crossing of t3 and t3b}]{(3,4,1)}{j31}
    \tnwire[kind=string, name=t3b]{j31}{j32}
    \tnwire[kind=string, name=t3c,
            cross={under at crossing of t3c and t3b}]{j32}{(3,4,2)}
    \tn[at=0.6 e of (4,4,1), skin=none, name=j41]{}
    \tn[at=0.6 e of (4,4,2), skin=none, name=j42]{}
    \tnwire[kind=string, name=t4, cross={under at crossing of t4 and t4b}]{(4,4,1)}{j41}
    \tnwire[kind=string, name=t4b]{j41}{j42}
    \tnwire[kind=string, name=t4c,
            cross={under at crossing of t4c and t4b}]{j42}{(4,4,2)}
  \end{tenkz} \end{gathered} \]
